\begin{abstract}
Visual place recognition (VPR) on the Moon has to cope with black shadows that move with the Sun,
statistically self-similar crater fields and strongly non-Lambertian regolith. We present
\emph{Galatea}, a framework that treats both stages of two-stage VPR as optimal transport (OT).
For global aggregation, a semi-relaxed unbalanced OT problem defines a one-parameter family of
token-to-cluster assignments that contains the NetVLAD ($\lambda=0$) and SALAD ($\lambda=1$)
assignments as end points and a learned $\lambda$ in between. We prove that its truncated Sinkhorn
iterations converge linearly at rate $\lambda$, characterise the dust-bin as a logistic gate, and
show that after intra-normalisation the descriptor sees the transport plan only through per-token
weights, which predicts when the choice of $\lambda$ can matter. For global--local fusion, a
training-free partial fused Gromov--Wasserstein matching scores both the appearance and the layout of
keypoints, and a saturating fusion rule provably overrides the global ranking only where the global
margin is small. On LunarSim-VPR, a physically motivated benchmark that renders the same query poses
under seven Sun conditions, terrestrial models fall from about $\zsSaladCzero\%$ to $\zsSaladHard\%$
Recall@1 under large Sun changes. With a frozen backbone, every value of $\lambda$ gives
$\approx\galRone\%$ (zero-shot: $\zsSaladRone\%$), as the theory predicts. Fine-tuning the last
backbone blocks adds $11$--$19$ points; within the OT family the assignments stay within about two points and a
learned $\lambda$ stays at its initial value, behaving like a fixed one, while NetVLAD-style residual
aggregation pulls ahead ($\ftNvRone\%$). Fusion adds
$2$--$7$ points for every stage-1 model and local matcher, because it keeps about half of the queries
that local evidence rescues while breaking few that the global ranking already had right; the full system
reaches $\msnvFgwRone\%$. Illumination-diverse training remains the largest single factor.
\end{abstract}

\begin{IEEEkeywords}
Localization, visual place recognition, space robotics, optimal transport, deep learning for visual perception.
\end{IEEEkeywords}
